% ============================================
% DOC-07
% Gerado de molambudos.md — build_miolo.py
% ============================================
\section*{\textbf{DOC-07} --- Diagnóstico Psiquiátrico (1979)}
\markboth{DOC-07 --- Diagnóstico Psiquiátrico (1979)}{}

\noindent \textit{DSM-II --- Dr. Heitor Oliveira}

\noindent \textbf{Paciente:} Joaquim Antônio Correia --- Registro 1.260
\textbf{Médico:} Dr. Heitor Oliveira, CRM-MG 4.892
\textbf{Data:} 14 de abril de 1979

\vspace{0.5em}\noindent\rule{\textwidth}{0.5pt}\vspace{0.5em}

\noindent \textbf{CID-8 / DSM-II}

\noindent \textbf{Eixo I --- Diagnóstico Principal:}
- \textbf{295.3 --- Esquizofrenia, Tipo Paranóide}
  - Justificativa: Delírios sistematizados de perseguição (acredita que uma \textquotedblleft{}criatura\textquotedblright{} o habita), alucinações visuais e táteis persistentes, discurso desconexo com neologismos (\textquotedblleft{}molambudo\textquotedblright{} como estado existencial), embotamento afetivo intercalado com explosões de agressividade.

\noindent \textbf{Eixo II --- Diagnóstico Secundário:}
- \textbf{301.7 --- Personalidade Psicopática}
  - Justificativa: Incapacidade de sentir remorso, conduta agressiva recorrente (três homicídios confirmados na juventude), manipulação de profissionais de saúde, ausência de vínculos afetivos.

\noindent \textbf{Eixo III --- Condições Médicas Associadas:}
- Icterícia crônica (olho direito --- possível sequela de desnutrição severa na infância, ou hepatopatia não tratada).
- Desnutrição proteico-calórica compatível com a cronologia ficcional: seca de 1915 e transferência em 1917.

\noindent \textbf{Eixo IV --- Gravidade dos Estressores Psicossociais:}
- Extrema.

\noindent \textbf{Eixo V --- Nível Mais Alto de Funcionamento Adaptativo no Último Ano:}
- 10/100 --- Perigo persistente para si e para terceiros.

\vspace{0.5em}\noindent\rule{\textwidth}{0.5pt}\vspace{0.5em}

\noindent \textbf{Observações Clínicas:}

\noindent O paciente foi admitido no Hospital Colônia em 1917, aos 10 anos de idade. Registros históricos indicam que chegou \textquotedblleft{}em estado de inanição, proveniente dos campos de concentração do Ceará, órfão de pai e mãe, com um olho amarelo e sem falar uma palavra que não fosse 'foooome'.\textquotedblright{}

\noindent Ao longo de 62 anos de institucionalização, o paciente apresentou:
- Episódios recorrentes de automutilação (arranhava a própria face até sangrar);
- Três homicídios confirmados (1922, 1937, 1948) --- todas as vítimas eram funcionários ou pacientes que tentaram se aproximar;
- Períodos de mutismo seletivo (recusava-se a falar com a equipe; chegou a ficar 7 anos sem dirigir a palavra a funcionários --- de 1949 a 1956);
- Redação compulsiva de um \textquotedblleft{}diário\textquotedblright{} que nunca permitiu que nenhum médico lesse integralmente;
- Fixação no tema da \textquotedblleft{}fome\textquotedblright{} e de uma \textquotedblleft{}criatura de olhos amarelos\textquotedblright{} que o acompanha.

\noindent O paciente demonstra inteligência preservada --- possivelmente acima da média --- mas inteiramente desviada para a construção de seu delírio. Sua capacidade de narrar sua própria história com coerência interna (dentro do sistema delirante) é notável.

\noindent \textbf{Prognóstico:} Reservado. Paciente irrecuperável. Manter em isolamento com sedação contínua.

\vspace{0.5em}\noindent\rule{\textwidth}{0.5pt}\vspace{0.5em}

\noindent \textit{Dr. Heitor Oliveira}
\textit{Psiquiatra Responsável --- Ala de Segurança Máxima}
\textit{Hospital Colônia de Barbacena}

\noindent \textit{P.S.: Dr. Heitor Oliveira desapareceu em 13 de junho de 1979, dois meses após assinar este laudo. O processo de desaparecimento nunca foi encerrado. Na ficha funcional do hospital, um colega anotou, à mão, na margem: \textquotedblleft{}Ele parou de comer três semanas antes. Dizia que não sentia mais fome --- e isso, disse ele, era o pior sintoma de todos.\textquotedblright{}}

\vspace{0.5em}
\noindent\textit{\textbf{↪ Rotas:} \rota{DOC-08} \rotasep \rota{DOC-04}}